% Research article: a single-column paper with an abstract, sections, a
% figure, a table, an equation and a bibliography.
%
% Recompile with the green button or Ctrl+Enter. Lines starting with % are
% comments: LaTeX ignores them.
\documentclass[11pt,a4paper]{article}

\usepackage[T1]{fontenc}
\usepackage[margin=2.5cm]{geometry}
\usepackage{amsmath,amssymb}   % equations and math symbols
\usepackage{graphicx}          % \includegraphics
\usepackage{booktabs}          % clean table rules: \toprule, \midrule, \bottomrule
\usepackage[hidelinks]{hyperref}        % clickable references and links

\title{A Clear and Specific Title for Your Article}
\author{First Author\thanks{Department, University. \texttt{first@example.edu}}
  \and Second Author\thanks{Institute, City.}}
\date{\today}

\begin{document}

\maketitle

\begin{abstract}
State the problem, the approach, the main result and why it matters, in
about 150 words. Write the abstract last: it is easier once the paper exists.
\end{abstract}

\section{Introduction}
\label{sec:intro}

Open with the question the paper answers. Cite earlier work with
\verb|\cite{key}|, for example the original \LaTeX{} manual~\cite{lamport1994}
and a classic paper~\cite{einstein1905}. Section~\ref{sec:method} describes
the method and Section~\ref{sec:results} the results.

\section{Method}
\label{sec:method}

Numbered equations get a label so the text can refer to them. Equation
\eqref{eq:gaussian} is the Gaussian integral:
\begin{equation}
  \int_{-\infty}^{\infty} e^{-x^{2}}\,dx = \sqrt{\pi}.
  \label{eq:gaussian}
\end{equation}

\section{Results}
\label{sec:results}

Table~\ref{tab:results} summarises the measurements and
Figure~\ref{fig:example} shows where a figure goes.

\begin{table}[htbp]
  \centering
  \caption{Example results. Captions go above tables.}
  \label{tab:results}
  \begin{tabular}{lcc}
    \toprule
    Sample & Value & Uncertainty \\
    \midrule
    A & 1.23 & 0.04 \\
    B & 2.34 & 0.05 \\
    C & 3.45 & 0.07 \\
    \bottomrule
  \end{tabular}
\end{table}

\begin{figure}[htbp]
  \centering
  % Drop an image onto the editor (or paste one) and Oberleaf writes the
  % \includegraphics line for you. Until then this box marks the place.
  \fbox{\parbox[c][4cm][c]{0.7\linewidth}{\centering Figure placeholder}}
  \caption{Example figure. Captions go below figures.}
  \label{fig:example}
\end{figure}

\section{Conclusion}

Say what was shown, what it means, and what remains open.

\bibliographystyle{plain}
\bibliography{references}

\end{document}
